% !TEX root = ../../ctfp-print.tex

\lettrine[lhang=0.17]{如}{果你向程序员提起单子}，很可能会谈到效应（effects）。对于数学家来说，单子则与代数有关。我们稍后会讨论代数——它们在编程中扮演重要角色——但首先我想让你直观了解它与单子的关系。目前这只是一个粗略的类比，请暂且接受这种表述方式。

代数涉及表达式的创建、操作和求值。表达式通过算符构建而成。考虑这个简单表达式：
$$x^2 + 2 x + 1$$
该表达式由变量（如$x$）、常量（如1或2）通过加法、乘法等算符组合而成。作为程序员，我们通常将表达式视为树状结构。

\begin{figure}[H]
  \centering
  \includegraphics[width=0.3\textwidth]{images/exptree.png}
\end{figure}

\noindent
树是一种容器，更一般地，表达式是存储变量的容器。在范畴论中，我们用自函子（endofunctor）表示容器。若将类型$a$赋给变量$x$，我们的表达式将具有类型$m\ a$，其中$m$是构建表达式树的自函子。（非平凡分支表达式通常通过递归定义的自函子创建。）

对表达式最常见的操作是什么？是替换：用表达式代替变量。例如在本例中，我们可以用$y - 1$替换$X$得到：
$$(y - 1)^2 + 2 (y - 1) + 1$$
具体过程如下：我们取一个类型为$m\ a$的表达式，并应用类型为$a \to m\ b$的变换（$b$代表$y$的类型），结果得到类型为$m\ b$的表达式。明确写出就是：
$$m\ a \to (a \to m\ b) \to m\ b$$
没错，这就是单子bind的签名。

以上是动机阐述。现在进入单子的数学定义。数学家使用不同于程序员的记号：他们习惯用字母$T$表示自函子，希腊字母$\mu$表示\code{join}，$\eta$表示\code{return}。这两个函数都是多态的，因此可以猜测它们对应自然变换。

因此在范畴论中，单子被定义为配备了一对自然变换$\mu$和$\eta$的自函子$T$。

$\mu$是从函子平方$T^2$到$T$的自然变换。这里的平方指函子与其自身的复合$T \circ T$（只有自函子才能进行此类"平方"运算）。
$$\mu \Colon T^2 \to T$$
该自然变换在对象$a$处的分量是态射：
$$\mu_a \Colon T (T a) \to T a$$
在$\Hask$范畴中，这直接对应我们熟悉的\code{join}定义。

$\eta$是恒等函子$I$到$T$的自然变换：
$$\eta \Colon I \to T$$
由于恒等函子作用于对象$a$的结果就是$a$本身，$\eta$在$a$处的分量由以下态射给出：
$$\eta_a \Colon a \to T a$$
这正好对应我们定义的\code{return}。

这些自然变换必须满足额外定律。一种理解方式是：这些定律允许我们为自函子$T$定义克莱斯利范畴（Kleisli category）。记住，克莱斯利箭头（arrow）$a \to b$定义为态射$a \to T b$。两个这样的箭头（用带下标$T$的圆圈表示）的复合可通过$\mu$实现：
$$g \circ_T f = \mu_c \circ (T g) \circ f$$
其中
\begin{gather*}
  f \Colon a \to T b \\
  g \Colon b \to T c
\end{gather*}
这里作为函子的$T$可以作用于态射$g$。用Haskell记号更容易识别这个公式：

\src{snippet01}
或者用分量形式表示：

\src{snippet02}
从代数解释的角度看，这只是连续两次替换的合成。

要使克莱斯利箭头构成范畴，要求其复合满足结合律，且$\eta_a$成为$a$处的单位箭头。这一要求可转化为$\mu$和$\eta$的单子定律。但另一种推导方式使其更像幺半群定律。事实上$\mu$常被称为\emph{乘法}，$\eta$被称为\emph{单位元}。

粗略地说，结合律要求以两种不同方式将$T$的立方$T^3$约简到$T$时结果相同。两个单位元律（左/右）要求当$\eta$作用于$T$后经$\mu$约简时，应回到原始的$T$。

此处有些微妙之处，因为我们是在组合自然变换和函子。因此需要复习水平组合（horizontal composition）知识。例如$T^3$可视作$T^2$后再作$T$的复合。我们可以对其应用两个自然变换的水平组合：
$$I_T \circ \mu$$

\begin{figure}[H]
  \centering
  \includegraphics[width=0.4\textwidth]{images/assoc1.png}
\end{figure}

\noindent
从而得到$T \circ T$；再通过$\mu$进一步约简为$T$。$I_T$是从$T$到自身的恒等自然变换。这类水平组合$I_T \circ \mu$的记号常简化为$T \circ \mu$。此记号无歧义，因为函子与自然变换直接组合没有意义，故此处$T$必然指$I_T$。

\noindent
我们也可以在自函子范畴${[}\cat{C}, \cat{C}{]}$中绘制该图示：

\begin{figure}[H]
  \centering
  \includegraphics[width=0.3\textwidth]{images/assoc2.png}
\end{figure}

\noindent
另一种方式是将$T^3$视作$T^2 \circ T$的复合并应用$\mu \circ T$。结果同样是$T \circ T$，再次通过$\mu$约简为$T$。我们要求这两种路径产生相同结果。

\begin{figure}[H]
  \centering
  \includegraphics[width=0.3\textwidth]{images/assoc.png}
\end{figure}

\noindent
类似地，我们可以将水平组合$\eta \circ T$应用于恒等函子$I$后再作$T$的复合，得到$T^2$，再通过$\mu$约简。结果应与直接对$T$应用恒等自然变换相同。同理，$T \circ \eta$的情况也应成立。

\begin{figure}[H]
  \centering
  \includegraphics[width=0.4\textwidth]{images/unitlawcomp-1.png}
\end{figure}

\noindent
你可以验证这些定律确实保证了克莱斯利箭头的复合满足范畴公理。

单子与幺半群的相似性令人印象深刻：都有乘法$\mu$、单位元$\eta$、结合律和单位律。但我们对幺半群的定义过于狭隘，不足以描述单子。因此需要推广幺半群的概念。

\section{幺半范畴}

回到传统的幺半群定义：它是一个带有二元运算和特殊单位元素的集合。在Haskell中可用类型类表达：

\src{snippet03}
二元运算\code{mappend}必须满足结合律和单位律（即与单位元\code{mempty}相乘不起作用）。

注意Haskell中\code{mappend}的定义是柯里化的。它可以解读为将每个\code{m}元素映射为函数：

\src{snippet04}
正是这种解读导致了将幺半群定义为单对象范畴的观点，其中自同态\code{(m -> m)}代表幺半群元素。但由于Haskell内置了柯里化，我们本可以从不同的乘法定义出发：

\src{snippet05}
此处笛卡尔积\code{(m, m)}成为待乘元素对的来源。

该定义暗示了另一种推广路径：用范畴积替代笛卡尔积。我们可以从一个全局定义了积的范畴出发，选取对象\code{m}，并定义乘法为态射：
$$\mu \Colon m\times{}m \to m$$
但存在一个问题：在任意范畴中无法观察对象内部结构，如何选取单位元素？这里有个技巧：记得元素选择等价于从单元素集的函数吗？在Haskell中可以用函数替代\code{mempty}的定义：

\src{snippet06}
单元素集是$\Set$范畴的终对象，因此可将此定义推广到任何具有终对象$t$的范畴：
$$\eta \Colon t \to m$$
这使我们无需谈论元素即可选取单位"元素"。

与之前将幺半群定义为单对象范畴不同，此时幺半群定律不会自动满足——必须显式施加。为此我们需要建立底层范畴积本身的幺半结构。先回顾Haskell中的幺半结构。

从结合律开始。Haskell对应的等式定律为：

\src{snippet07}
要推广到其他范畴，需要将其重写为函数（态射）相等式。必须抽象掉对个别变量的作用——换句话说，使用无点（point-free）记法。利用笛卡尔积作为双函子的性质，左边可写为：

\src{snippet08}
右边可写为：

\src{snippet09}
这几乎达到目的。遗憾的是笛卡尔积并非严格结合——\code{(x, (y, z))}不等于\code{((x, y), z)}——因此不能直接写成无点形式：

\src{snippet10}
另一方面，这两种配对嵌套方式是同构的。存在称为结合子（associator）的可逆函数在两者间转换：

\src{snippet11}
借助结合子，可以写出\code{mu}的无点结合律：

\src{snippet12}
对单位律也可采用类似技巧，新记法下的定律形式为：

\src{snippet13}
可重写为：

\src{snippet14}
这些同构\code{lambda}和\code{rho}分别称为左/右单位子。它们表明单位元\code{()}按同构是笛卡尔积的单位：

\src{snippet15}

\src{snippet16}
因此单位律的无点版本为：

\src{snippet17}
我们已利用底层笛卡尔积在类型范畴中表现得像幺半群乘法的事实，建立了\code{mu}和\code{eta}的无点幺半律。但请注意，笛卡尔积的结合律和单位律仅在同构意义下成立。

事实证明这些定律可推广到任何具有积和终对象的范畴。范畴积确实按同构满足结合律，终对象作为单位元也是如此。结合子和两个单位子都是自然同构。这些定律可用交换图表示。

\begin{figure}[H]
  \centering
  \includegraphics[width=0.5\textwidth]{images/assocmon.png}
\end{figure}

\noindent
注意由于积是双函子，它可以提升一对态射——在Haskell中这是通过\code{bimap}实现的。

我们可以在此止步，声明在任何具有范畴积和终对象的范畴上都能定义幺半群。只要能选取对象$m$及满足幺半律的两个态射$\mu$和$\eta$，就构成了幺半群。但我们还能做得更好：定义$\mu$和$\eta$的定律并不需要完整的范畴积。回想积通过使用投影的泛构造定义，而在幺半律的表述中并未使用任何投影。

行为类似积但不是积的双函子称为\newterm{张量积}（tensor product），常用中缀运算符$\otimes$表示。一般而言张量积的定义较为复杂，但无需深究，只需列举其性质——最重要的是按同构满足结合律。

同样地，对象$t$不必是终对象。我们从未使用其终性质——即从任何对象到它的唯一态射的存在性。需要的是它能与张量积良好配合，即要求它是张量积的单位元（同样按同构）。综合起来：

幺半范畴是配备了一个称为张量积的双函子：
$$\otimes \Colon \cat{C}\times{}\cat{C} \to \cat{C}$$
以及一个特殊的单位对象$i$，并带有三个自然同构分别称为结合子和左/右单位子：
\begin{align*}
  \alpha_{a b c} & \Colon (a \otimes b) \otimes c \to a \otimes (b \otimes c) \\
  \lambda_a      & \Colon i \otimes a \to a                                   \\
  \rho_a         & \Colon a \otimes i \to a
\end{align*}
（还有一个简化四重张量积的协调条件。）

关键在于张量积描述了许多常见的双函子。尤其适用于积、余积，以及即将看到的自函子复合（还有诸如Day卷积等更晦涩的积）。幺半范畴将在富范畴（enriched categories）的表述中发挥核心作用。

\section{幺半范畴中的幺半群}

我们现在可以在更一般的幺半范畴(monoidal category)框架下定义幺半群(monoid)了。我们首先选取一个对象 $m$。通过张量积(tensor product)，我们可以构造 $m$ 的各次幂：$m$ 的平方是 $m \otimes m$。构造 $m$ 的立方有两种方式，但它们通过结合子(associator)是同构的。类似地，$m$ 的更高次幂也成立（这就是我们需要协调条件(coherence conditions)的原因）。要构造一个幺半群，我们需要选取两个态射：
\begin{align*}
  \mu  & \Colon m \otimes m \to m \\
  \eta & \Colon i \to m
\end{align*}
其中 $i$ 是张量积的单位对象(unit object)。

\begin{figure}[H]
  \centering
  \includegraphics[width=0.4\textwidth]{images/monoid-1.jpg}
\end{figure}

\noindent
这些态射必须满足结合律和单位律，可以用以下交换图(commuting diagram)来表述：

\begin{figure}[H]
  \centering
  \includegraphics[width=0.5\textwidth]{images/assoctensor.jpg}
\end{figure}

\begin{figure}[H]
  \centering
  \includegraphics[width=0.5\textwidth]{images/unitmon.jpg}
\end{figure}

\noindent
需要注意的是，张量积必须是一个双函子(bifunctor)这一点至关重要，因为我们需要提升态射对来构造诸如 $\mu \otimes \id$ 或 $\eta \otimes \id$ 这样的乘积。这些交换图只是我们之前关于范畴积结果的直接推广。

\section{作为幺半群的单子}

张量积结构(monoidal structures)常常出现在意想不到的地方。函子范畴(functor category)就是这样一个场所。如果你稍加观察，或许能发现函子复合可以视为某种乘法形式。问题在于并非任意两个函子都能复合——其中一个的靶范畴必须等于另一个的源范畴。这正是态射复合的常规法则，正如我们所知，函子确实是范畴 $\Cat$ 中的态射。但就像自同态(endomorphisms)（返回同一对象的态射）总是可复合的，自函子(endofunctors)也是如此。对于给定范畴 $\cat{C}$，从 $\cat{C}$ 到自身的自函子构成函子范畴 ${[}\cat{C}, \cat{C}{]}$。其对象是自函子，态射是它们之间的自然变换(natural transformations)。我们可以从该范畴中任取两个对象，设为自函子 $F$ 和 $G$，并构造第三个对象 $F \circ G$——它们的复合自函子。

自函子复合能否成为张量积的良好候选？首先需验证它是否构成双函子(bifunctor)。它能否用于提升一对态射——在此即自然变换？张量积对应的 \code{bimap} 签名应形如：
$$\mathit{bimap} \Colon (a \to b) \to (c \to d) \to (a \otimes c \to b \otimes d)$$
若将对象替换为自函子，箭头替换为自然变换，张量积替换为复合，则得到：
$$(F \to F') \to (G \to G') \to (F \circ G \to F' \circ G')$$
这正是水平复合(horizontal composition)的特例。

\begin{figure}[H]
  \centering
  \includegraphics[width=0.4\textwidth]{images/horizcomp.png}
\end{figure}

\noindent
我们还拥有恒等自函子 $I$，它可充当自函子复合的单位元——我们的新张量积。更重要的是，函子复合具有结合性。事实上，此处的结合律和单位律是严格成立的——无需引入结合约束(associator)或两个单位约束(unitors)。因此自函子范畴以函子复合为张量积构成严格的单子范畴(strict monoidal category)。

此范畴中的幺半群(monoid)是什么？它包含一个对象——即自函子 $T$；以及两个态射——即自然变换：
\begin{gather*}
  \mu \Colon T \circ T \to T \\
  \eta \Colon I \to T
\end{gather*}
不仅如此，这些幺半群律(moid laws)如下：

\begin{figure}[H]
  \centering
  \includegraphics[width=0.4\textwidth]{images/assoc.png}
\end{figure}

\begin{figure}[H]
  \centering
  \includegraphics[width=0.5\textwidth]{images/unitlawcomp.png}
\end{figure}

\noindent
这正是我们之前见过的单子律(monad laws)。现在你应当理解Saunders Mac Lane的著名论断了：

\begin{quote}
  总而言之，单子不过是自函子范畴中的幺半群。
\end{quote}
你或许曾在函数式编程会议的某些T恤上见过这句话大放异彩。

\section{从伴随关系生成的单子}

一个\hyperref[adjunctions]{伴随}\footnote{参见第18章关于\hyperref[adjunctions]{伴随}的内容}
$L \dashv R$，是由两个范畴$\cat{C}$和$\cat{D}$之间往返的一对函子构成。
通过两种组合方式可得到两个自函子(endofunctor)：$R \circ L$和$L \circ R$。
根据伴随关系的定义，这两个自函子通过称为单位元(unit)和余单位元(counit)的两个自然变换与恒等函子相关联：
\begin{gather*}
  \eta \Colon I_{\cat{D}} \to R \circ L \\
  \varepsilon \Colon L \circ R \to I_{\cat{C}}
\end{gather*}
立即可见伴随关系的单位元与单子的单位元形式相同。事实上自函子$R \circ L$确实构成单子。
我们只需定义合适的$\mu$来配合$\eta$即可。这是自函子平方与其自身的自然变换，用伴随函子表示即：
$$R \circ L \circ R \circ L \to R \circ L$$
实际上我们可以利用余单位元坍缩中间的$L \circ R$。$\mu$的具体公式由水平复合给出：
$$\mu = R \circ \varepsilon \circ L$$
单子公理(monadic laws)可从伴随关系的单位元和余单位元满足的恒等式及交换律推导得出。

在Haskell中我们很少见到从伴随关系派生的单子，因为伴随关系通常涉及两个不同范畴。然而指数对象(exponential/object)或函数对象(function object)的定义是个例外。以下是形成该伴随关系的两个自函子：
\begin{gather*}
  L z = z\times{}s \\
  R b = s \Rightarrow b
\end{gather*}
您可能已认出它们的组合就是熟悉的state单子：
$$R (L z) = s \Rightarrow (z\times{}s)$$
我们在Haskell中见过这个单子：

\src{snippet18}
现在也将伴随关系翻译到Haskell中。左函子是乘积函子：

\src{snippet19}
右函子是读者函子(reader functor)：

\src{snippet20}
它们构成伴随关系：

\src{snippet21}
您可以验证读者函子与乘积函子的组合确实等价于状态函子：

\src{snippet22}
正如预期，伴随关系的\code{unit}对应state单子的\code{return}函数。\code{counit}的作用是对函数作用于其参数进行求值。这正是\code{runState}函数的非柯里化版本：

\src{snippet23}
(非柯里化是因为\code{counit}作用于序对)

现在我们可以将\code{join}定义为自然变换$\mu$的分量。为此需要三个自然变换的水平复合：
$$\mu = R \circ \varepsilon \circ L$$
换句话说我们需要让余单位元$\varepsilon$穿过一层读者函子(reader functor)。不能直接调用\code{fmap}，因为编译器会选取\code{State}函子而非\code{Reader}函子的实现。但要注意读者函子的\code{fmap}本质上是左函数组合。因此我们将直接使用函数组合操作。

首先需要剥离\code{State}数据构造器以暴露内部函数，这通过\code{runState}实现：

\src{snippet24}
然后将其左组合(left-compose)余单位元，后者定义为\code{uncurry runState}。最后重新包装进\code{State}数据构造器：

\src{snippet25}
这正是\code{State}单子\code{join}的标准实现。

值得注意的是不仅每个伴随关系都能生成单子，反过来也成立：每个单子都可以分解为两个伴随函子的组合。不过这种分解并非唯一。

我们将在下一节讨论另一个自函子$L \circ R$。